% =================================================================
%  CSE 815 -- Machine Learning
%  Module 1, Week 2, Part 2
%  Feature Scaling, Convergence, Learning Rate, Feature Engineering
%  and Polynomial Regression
% =================================================================
\documentclass[aspectratio=169,10pt]{beamer}
\usepackage{cse815}

\title[CSE 815 -- Week 2, Part 2]{Making Gradient Descent Work in Practice}
\subtitle{Feature Scaling, Convergence, the Learning Rate, and Feature Engineering}
\author{Rokan Uddin Faruqui\\ \small Professor, Department of CSE}
\institute{\small University of Chittagong \\ \texttt{rufaruqui@cu.ac.bd}}
\date[Week 2 / Part 2]{Week 2 $\cdot$ Part 2}

\begin{document}

\begin{frame}[plain,noframenumbering]
  \titlepage
\end{frame}

% ---------------------------------------------------------------
\begin{frame}{From correct to usable}
  \begin{columns}[T,onlytextwidth]
    \begin{column}{0.5\textwidth}
      \begin{block}{Part 1 gave us}
        \[ f_{\vw,b}(\vx)=\vw\cdot\vx+b \]
        \[ \vw := \vw - \alpha\,\nabla_{\vw}J,\quad b := b - \alpha\,\tfrac{\partial J}{\partial b} \]
      \end{block}
      \footnotesize
      That algorithm is correct. Run it on the flats table as written and it will crawl,
      or diverge. Today is about the difference between an algorithm that is right and one
      that actually finishes.
    \end{column}
    \begin{column}{0.46\textwidth}
      \begin{keyidea}[Today]
        \footnotesize
        \begin{enumerate}
          \item Why unscaled features cripple descent %\clo{15 min}
          \item Feature scaling: three methods %\clo{20 min}
          \item Checking for convergence %\clo{10 min}
          \item Choosing $\alpha$ systematically %\clo{15 min}
          \item Feature engineering %\clo{10 min}
          \item Polynomial regression %\clo{15 min}
        \end{enumerate}
      \end{keyidea}
    \end{column}
  \end{columns}
\end{frame}

% ===============================================================
\section{Feature Scaling}

\begin{frame}{The problem: features on wildly different scales}
  \footnotesize
  Take two features of a flat: size $x_1 \in [300, 2000]$ sq ft, and bedrooms
  $x_2 \in [0,5]$. A flat at $x_1=2000,\ x_2=5$ sells for 50 lakh.
  \vspace{1mm}
  \begin{columns}[T,onlytextwidth]
    \begin{column}{0.48\textwidth}
      \begin{block}{Which parameters are plausible?}
        $w_1=50,\ w_2=0.1,\ b=50$ gives
        \[ 50(2000)+0.1(5)+50 \approx 100{,}050 \quad\text{\alert{far too large}} \]
        $w_1=0.1,\ w_2=50,\ b=50$ gives
        \[ 0.1(2000)+50(5)+50 = 500 \quad\text{\alert{still too large}} \]
        Only a \emph{tiny} $w_1$ can work, because $x_1$ is huge.
      \end{block}
    \end{column}
    \begin{column}{0.48\textwidth}
      \begin{keyidea}
        A feature with a \textbf{large} range needs a \textbf{small} weight, and vice versa.
        So the cost surface is stretched: $J$ is hypersensitive to $w_1$ and nearly flat
        in $w_2$.
      \end{keyidea}
    \end{column}
  \end{columns}
\end{frame}

\begin{frame}{What that does to the contours}
  \begin{columns}[T,onlytextwidth]
    \begin{column}{0.5\textwidth}
      \centering {\scriptsize Unscaled: a long, thin, tilted valley}\par
      \begin{tikzpicture}
        \begin{axis}[width=6.0cm,height=4.4cm,xlabel={$w_1$},ylabel={$w_2$},
          tick label style={font=\scriptsize},label style={font=\scriptsize},
          view={0}{90}, domain=-1:1, y domain=-9:9,
          colormap={cu}{color=(cublue) color=(cuteal) color=(cuamber)}]
          \costcontour{}{0.4,1.2,3,7,14,26}{40*x^2 + 0.28*y^2}
          \addplot3[curose,thick,mark=*,mark size=1.2pt] coordinates {
            (-0.72,-7.6,0) (0.30,-7.0,0) (-0.13,-6.5,0) (0.05,-6.0,0)
            (-0.02,-5.6,0) (0.01,-5.2,0) (0.0,-4.8,0)};
          \addplot3[only marks,mark=star,mark size=3pt,cublue] coordinates {(0,0,0)};
        \end{axis}
      \end{tikzpicture}
      {\scriptsize Descent oscillates across the valley.}
    \end{column}
    \begin{column}{0.48\textwidth}
      \centering {\scriptsize Scaled: near-circular contours}\par
      \begin{tikzpicture}
        \begin{axis}[width=6.0cm,height=4.4cm,xlabel={$w_1$},ylabel={$w_2$},
          tick label style={font=\scriptsize},label style={font=\scriptsize},
          view={0}{90}, domain=-1:1, y domain=-1:1,
          colormap={cu}{color=(cublue) color=(cuteal) color=(cuamber)}]
          \costcontour{}{0.08,0.2,0.42,0.72,1.1,1.6}{2.1*x^2 + 1.8*y^2}
          \addplot3[curose,thick,mark=*,mark size=1.2pt] coordinates {
            (-0.82,-0.78,0) (-0.42,-0.42,0) (-0.20,-0.22,0) (-0.09,-0.11,0) (-0.03,-0.05,0)};
          \addplot3[only marks,mark=star,mark size=3pt,cublue] coordinates {(0,0,0)};
        \end{axis}
      \end{tikzpicture}
      {\scriptsize Descent heads almost straight for the minimum.}
    \end{column}
  \end{columns}
  \vspace{1mm}
  \begin{keyidea}
    \footnotesize
    Steepest descent is perpendicular to the local contour. In a thin valley that direction
    points \emph{across} the valley, not along it, so the path zig-zags and $\alpha$ must stay
    tiny to avoid overshooting. Round contours remove the problem.
  \end{keyidea}
\end{frame}

\begin{frame}{Three ways to scale}
  \footnotesize
  Throughout, $x_1$ ranges over $[300, 2000]$ with mean $\mu_1 = 600$ and standard
  deviation $\sigma_1 = 450$.
  \vspace{1mm}
  \begin{columns}[T,onlytextwidth]
    \begin{column}{0.32\textwidth}
      \begin{block}{Divide by the max}
        \[ x_{1,\text{scaled}} = \frac{x_1}{\max} \]
        \[ \frac{300}{2000} \le x_1 \le \frac{2000}{2000} \]
        Gives $0.15 \le x_1 \le 1$.\\[2pt]
        Simplest; range is not centred.
      \end{block}
    \end{column}
    \begin{column}{0.32\textwidth}
      \begin{block}{Mean normalization}
        \[ x_{1,\text{scaled}} = \frac{x_1-\mu_1}{\max-\min} \]
        \[ \frac{300-600}{1700} \le \cdot \le \frac{2000-600}{1700} \]
        Gives $-0.18 \le x_1 \le 0.82$.\\[2pt]
        Centred on zero.
      \end{block}
    \end{column}
    \begin{column}{0.32\textwidth}
      \begin{block}{Z-score normalization}
        \[ x_{1,\text{scaled}} = \frac{x_1-\mu_1}{\sigma_1} \]
        \[ \frac{300-600}{450} \le \cdot \le \frac{2000-600}{450} \]
        Gives $-0.67 \le x_1 \le 3.11$.\\[2pt]
        The usual default.
      \end{block}
    \end{column}
  \end{columns}
  \vspace{2mm}
  \begin{keyidea}
    Aim for roughly $-1 \le x_j \le 1$ for every feature. The bound is a rule of thumb, not a
    law: $-3\le x_j\le 3$ is fine, $-0.3\le x_j\le 0.3$ is fine. Rescale when one feature's
    range is orders of magnitude away from the others.
  \end{keyidea}
\end{frame}

\begin{frame}{Scaling in practice --- two rules that matter}
  \footnotesize
  \begin{columns}[T,onlytextwidth]
    \begin{column}{0.48\textwidth}
      \begin{pitfall}[Fit on train, apply everywhere]
        Compute $\mu_j$ and $\sigma_j$ from the \emph{training set only}, store them, and
        reuse those same numbers for validation, test and production data. Recomputing them
        on the test set leaks information and quietly inflates your reported score.
      \end{pitfall}
    \end{column}
    \begin{column}{0.48\textwidth}
      \begin{pitfall}[Scaling changes what $\vw$ means]
        After scaling, $w_j$ is the effect of a \emph{one-standard-deviation} change in
        feature $j$, not of one square foot. To quote a coefficient in original units you
        must unscale it. Predictions are unaffected; interpretations are not.
      \end{pitfall}
    \end{column}
  \end{columns}
  \vspace{1mm}
  \begin{block}{In code}
    \texttt{mu\ \ \ \ = X\_train.mean(axis=0)\ \ \ \ \ \# per-feature, shape (n,)}\\
    \texttt{sigma\ = X\_train.std(axis=0)}\\
    \texttt{X\_train\_s = (X\_train - mu) / sigma}\\
    \texttt{X\_test\_s\ = (X\_test\ - mu) / sigma\ \ \ \# same mu, sigma}
  \end{block}
  \small
  \texttt{axis=0} collapses rows, giving one statistic per feature --- the single most common
  place to get this wrong.
\end{frame}

% ===============================================================
\section{Convergence and the Learning Rate}

\begin{frame}{Checking gradient descent for convergence}
  \begin{columns}[T,onlytextwidth]
    \begin{column}{0.52\textwidth}
      \centering {\scriptsize The learning curve is the diagnostic}\par
      \begin{tikzpicture}
        \begin{axis}[width=6.4cm,height=4.4cm,xlabel={iteration},ylabel={$J(\vw,b)$},
          xmin=0,xmax=400,ymin=0,ymax=1.05,grid=major,grid style={culight},
          tick label style={font=\scriptsize},label style={font=\scriptsize}]
          \addplot[cublue,very thick,domain=0:400,samples=120]{exp(-0.016*x)+0.02};
          \draw[dashed,cugray] (axis cs:0,0.05)--(axis cs:400,0.05);
          \node[font=\scriptsize,cugray,anchor=south east] at (axis cs:390,0.09) {converged};
          \draw[-{Stealth},cuteal,thick] (axis cs:80,0.55)--(axis cs:110,0.30);
          \node[font=\scriptsize,cuteal,anchor=west] at (axis cs:112,0.34) {still improving};
        \end{axis}
      \end{tikzpicture}
    \end{column}
    \begin{column}{0.44\textwidth}
      \footnotesize
      \begin{itemize}
        \item Plot $J$ against \emph{iteration number}, not against a parameter.
        \item $J$ must decrease on \textbf{every} iteration. One rise means $\alpha$ is too
              large or the gradient is wrong.
        \item The curve flattening is convergence.
        \item The number of iterations needed varies enormously --- 30 for one problem,
              100{,}000 for another. There is no way to know in advance; you look.
      \end{itemize}
      \begin{defnbox}[Automatic convergence test]
        \footnotesize
        Declare convergence when $J$ falls by less than $\varepsilon$ (say $10^{-3}$) in one
        iteration. Convenient, but choosing $\varepsilon$ is itself hard --- which is why the
        plot is preferred.
      \end{defnbox}
    \end{column}
  \end{columns}
\end{frame}

\begin{frame}{Choosing the learning rate}
  \footnotesize
  \begin{columns}[T,onlytextwidth]
    \begin{column}{0.5\textwidth}
      \centering {\scriptsize Three ways a run can look}\par
      \begin{tikzpicture}
        \begin{axis}[width=6.2cm,height=4.2cm,xlabel={iteration},ylabel={$J$},
          xmin=0,xmax=40,ymin=0,ymax=2.2,grid=major,grid style={culight},
          tick label style={font=\scriptsize},label style={font=\scriptsize},
          legend style={font=\tiny,at={(0.98,0.98)},anchor=north east,draw=none,fill=white}]
          \addplot[cuteal,very thick,domain=0:40,samples=60]{1.6*exp(-0.22*x)+0.05};
          \addlegendentry{good}
          \addplot[cublue,very thick,dashed,domain=0:40,samples=60]{1.6*exp(-0.02*x)+0.05};
          \addlegendentry{$\alpha$ too small}
          \addplot[curose,very thick,domain=0:40,samples=60]{0.35+0.30*exp(0.045*x)*abs(sin(deg(0.9*x)))};
          \addlegendentry{$\alpha$ too large}
        \end{axis}
      \end{tikzpicture}
    \end{column}
    \begin{column}{0.46\textwidth}
      \begin{block}{The procedure}
        Try a range of values, each about $3\times$ the last:
        \[ 0.001,\ 0.003,\ 0.01,\ 0.03,\ 0.1,\ 0.3,\ 1 \]
        Run each for a small, fixed number of iterations and plot $J$. Pick the
        \emph{largest} $\alpha$ whose curve still decreases smoothly, then back off one step.
      \end{block}
      \begin{pitfall}
        If $J$ rises with a very small $\alpha$ --- say $10^{-6}$ --- the bug is in the
        gradient, not the learning rate. A correct gradient with a tiny step \emph{must}
        decrease $J$.
      \end{pitfall}
    \end{column}
  \end{columns}
\end{frame}

% ===============================================================
\section{Feature Engineering and Polynomial Regression}

\begin{frame}{Feature engineering}
  \footnotesize
  \begin{columns}[T,onlytextwidth]
    \begin{column}{0.5\textwidth}
      Predicting the price of a plot of land from its frontage $x_1$ and depth $x_2$:
      \[ f(\vx) = w_1x_1 + w_2x_2 + b . \]
      This treats frontage and depth as independent contributions. But what actually drives
      price is \emph{area}:
      \[ x_3 = x_1 x_2, \qquad f(\vx) = w_1x_1 + w_2x_2 + w_3x_3 + b . \]
      The model can now express something no combination of $w_1,w_2$ could.
      \vspace{2mm}
      \begin{defnbox}[Feature engineering]
        Using domain knowledge to build new features from existing ones --- by transforming
        or combining them --- so that the model can represent the relationship you believe
        is there.
      \end{defnbox}
    \end{column}
    \begin{column}{0.46\textwidth}
      \centering
      \begin{tikzpicture}[font=\scriptsize]
        \draw[fill=culight,draw=cublue,thick] (0,0) rectangle (3.4,2.0);
        \draw[<->,cuamber,thick] (0,-0.28)--(3.4,-0.28)
          node[midway,below,cuamber]{frontage $x_1$};
        \draw[<->,cuteal,thick] (3.66,0)--(3.66,2.0)
          node[midway,right,cuteal]{depth $x_2$};
        \node[cublue] at (1.7,1.0) {area $x_3 = x_1x_2$};
      \end{tikzpicture}
      \vspace{3mm}
      \begin{keyidea}
        \footnotesize
        The choice of features often matters more than the choice of algorithm. A linear
        model on well-chosen features routinely beats a complicated model on raw ones.
      \end{keyidea}
    \end{column}
  \end{columns}
\end{frame}

\begin{frame}{Polynomial regression}
  \begin{columns}[T,onlytextwidth]
    \begin{column}{0.5\textwidth}
      \centering
      \begin{tikzpicture}
        \begin{axis}[width=6.4cm,height=4.6cm,
          xlabel={size $x$ (k sq ft)},ylabel={price (lakh)},
          xmin=0,xmax=4.4,ymin=0,ymax=95,grid=major,grid style={culight},
          tick label style={font=\scriptsize},label style={font=\scriptsize},
          legend style={font=\tiny,at={(0.98,0.03)},anchor=south east,draw=none,fill=white}]
          \addplot[only marks,mark=*,mark size=1.7pt,cublue,forget plot]
            coordinates {(0.5,18)(0.9,31)(1.3,43)(1.7,52)(2.1,60)(2.5,66)(2.9,71)(3.3,74)(3.8,77)};
          \addplot[curose,very thick,domain=0.3:4.2]{18.6*x+11.5};
          \addlegendentry{linear}
          \addplot[cuteal,very thick,domain=0.3:4.2]{-4.6*x^2+37*x+2.5};
          \addlegendentry{quadratic}
          \addplot[cuamber,very thick,dashed,domain=0.3:4.2]{40*sqrt(x)-8};
          \addlegendentry{square root}
        \end{axis}
      \end{tikzpicture}
    \end{column}
    \begin{column}{0.46\textwidth}
      \footnotesize
      Price rises with size but \emph{flattens}. A straight line cannot bend. Feed the model
      powers of $x$ as extra features:
      \[ f(x) = w_1x + w_2x^2 + w_3x^3 + b \]
      \begin{itemize}
        \item Still \textbf{linear regression} --- linear in the \emph{parameters}. Only the
              features changed.
        \item A quadratic eventually turns downward; prices that only flatten are often
              better served by $\sqrt{x}$, which rises forever but ever more slowly.
      \end{itemize}
      \begin{pitfall}
        \footnotesize
        Powers explode the scale: if $x\in[1,4]$ then $x^3\in[1,64]$. \textbf{Feature scaling
        becomes mandatory} once you add polynomial features.
      \end{pitfall}
    \end{column}
  \end{columns}
\end{frame}

\begin{frame}{Choosing among candidate feature sets}
  \footnotesize
  You now have a way to generate endless models: $x$, $x^2$, $\sqrt{x}$, $x_1x_2$, and any
  combination. Which do you pick?
  \vspace{1mm}
  \begin{columns}[T,onlytextwidth]
    \begin{column}{0.48\textwidth}
      \begin{block}{What you can do this week}
        \begin{itemize}
          \item Plot the data and look at the shape.
          \item Use domain knowledge --- area is a real thing; $x^{7}$ is not.
          \item Compare the fitted $w_j$: a feature whose weight comes out near zero is not
                earning its place.
        \end{itemize}
      \end{block}
    \end{column}
    \begin{column}{0.48\textwidth}
      \begin{alertblock}{What you cannot do yet}
        Judging a model by its cost on the \emph{training} data always favours more
        features. Adding $x^{10}$ can drive $J$ toward zero while making predictions worse.
        \vspace{1mm}
        The tools for this --- train/validation/test splits, bias and variance,
        regularization --- arrive in Week 3 and Module 2.
      \end{alertblock}
    \end{column}
  \end{columns}
  \vspace{1mm}
  \begin{keyidea}
    Treat every engineered feature as a hypothesis about the world, and be ready to
    discard it.
  \end{keyidea}
\end{frame}

% ===============================================================
\section{Wrap-up and Assessment}

\begin{frame}{Summary of Week 2}
  \footnotesize
  \begin{columns}[T,onlytextwidth]
    \begin{column}{0.5\textwidth}
      \begin{block}{Concepts}
        \begin{itemize}
          \item Features on different scales stretch the cost contours, so descent zig-zags.
          \item Z-score normalization is the usual fix; fit $\mu,\sigma$ on the training set
                and reuse them.
          \item Plot $J$ against iteration. It must decrease, always.
          \item Choose $\alpha$ by trying $\times 3$ steps and taking the largest that still
                decreases smoothly.
          \item New features are hypotheses; polynomial terms are just more features.
        \end{itemize}
      \end{block}
    \end{column}
    \begin{column}{0.48\textwidth}
      \begin{block}{Formulas}
        \[ x_{j,\text{scaled}} = \frac{x_j-\mu_j}{\sigma_j} \]
        \[ f(x) = w_1x + w_2x^2 + b \quad\text{(still linear regression)} \]
      \end{block}
      \begin{block}{Next week}
        \footnotesize
        Classification: the logistic regression model, decision boundaries, a cost function
        that squared error cannot provide, and the problem of overfitting.
      \end{block}
    \end{column}
  \end{columns}
\end{frame}

\begin{frame}{Descriptive questions}
  \footnotesize
  \begin{enumerate}
    \item Explain why a feature with a large numerical range tends to receive a small weight,
          and what that does to the shape of the cost contours. \hfill\clo{6}
    \item Define mean normalization and z-score normalization, giving the formula for each.
          For $x\in[300,2000]$ with $\mu=600$, $\sigma=450$, compute the scaled range under
          both. \hfill\clo{7}
    \item Why must $\mu_j$ and $\sigma_j$ be computed on the training set alone? Describe
          concretely what goes wrong if they are recomputed on the test set. \hfill\clo{6}
    \item Describe the learning-curve test for convergence and state the one property the
          curve must have. Contrast it with an automatic $\varepsilon$-threshold test.
          \hfill\clo{6}
    \item Give the procedure for choosing $\alpha$. Explain why the values are spaced by
          roughly $3\times$ rather than uniformly. \hfill\clo{5}
  \end{enumerate}
\end{frame}

\begin{frame}{Descriptive questions (continued)}
  \footnotesize
  \begin{enumerate}\setcounter{enumi}{5}
    \item If $J$ increases while $\alpha = 10^{-6}$, what do you conclude and why? What is
          your next debugging step? \hfill\clo{5}
    \item Define feature engineering. Give an example, other than area, where a product or
          ratio of two raw features is more meaningful than either alone. \hfill\clo{5}
    \item ``Polynomial regression is not linear regression.'' Evaluate this claim and justify
          your verdict precisely. \hfill\clo{5}
    \item Explain why polynomial features make feature scaling mandatory. Illustrate with
          $x\in[1,4]$ and the features $x,x^2,x^3$. \hfill\clo{5}
    \item Why can a model with more features always achieve a lower training cost, and why is
          that not evidence that it is better? \hfill\clo{5}
  \end{enumerate}
\end{frame}

\begin{frame}{Design questions}
  \footnotesize
  \begin{enumerate}
    \item \textbf{Design a scaling pipeline.} Specify a reusable component that scales
          features for training, validation and production. State what it stores, its
          interface, and how it behaves when production data contains a value far outside
          the training range. \hfill\clo{12}
    \item \textbf{Design an $\alpha$ search.} Extend the $\times 3$ procedure into an
          automated routine that returns a recommended $\alpha$ without a human looking at
          plots. Define the decision rule, what you log, and one failure case your design
          does not cover. \hfill\clo{12}
    \item \textbf{Design the features.} You must predict daily rice-market prices in a
          Chattogram upazila. Propose six features, at least two of which are engineered
          from raw measurements, and justify each as a hypothesis about the market. Identify
          any feature that risks leaking the target. \hfill\clo{12}
    \item \textbf{Design a diagnosis procedure.} A colleague reports that training ``does not
          work'' on a 40-feature data set. Design an ordered checklist that distinguishes
          unscaled features, a bad $\alpha$, and a wrong gradient. Each step must state what
          you run and what result implicates which cause. \hfill\clo{12}
    \item \textbf{Choose a functional form.} Given data where price rises steeply then almost
          plateaus, compare a quadratic against $\sqrt{x}$ as the engineered feature. Argue
          from the behaviour of each function outside the training range which you would
          deploy. \hfill\clo{10}
  \end{enumerate}
\end{frame}

\begin{frame}{Reading}
  \begin{block}{Primary text}
    G.~James, D.~Witten, T.~Hastie and R.~Tibshirani,
    \emph{An Introduction to Statistical Learning}, 2nd ed., Springer, 2021.\\[2pt]
    \hfill{\small$\rightarrow$ \textbf{\S3.3} --- extensions of the linear model: interaction
    terms (feature engineering) and non-linear transformations of the predictors
    (polynomial regression).}
  \end{block}
  \begin{block}{Supplementary}
    C.~M.~Bishop, \emph{Pattern Recognition and Machine Learning}, Springer, 2006.\\[2pt]
    \hfill{\small$\rightarrow$ \textbf{\S1.1} --- polynomial curve fitting, and the first
    warning about model complexity that Week 3 takes up as overfitting.}
  \end{block}
  \vspace{2mm}
  \begin{center}
    \small\color{cugray}
    Lab (CSE 816, Week 2): feature scaling and the learning rate; feature engineering and
    polynomial regression; linear regression with scikit-learn.\\
    Practice quiz: \emph{Gradient descent in practice}.
  \end{center}
\end{frame}

\end{document}
